% ---------------------------------------------------------------------------
% Results and Discussion (Section 5) - drafted for M4.
% Numbers come from notebooks/06_results_error_analysis.ipynb (Parts 1-3) and
% results/metrics/{test_results,test_significance,efficiency,data_size_curve}.csv.
% The main table is generated: results/tables/test_results.tex.
% ---------------------------------------------------------------------------
\section{Results and Discussion}
\label{sec:results}

\subsection{How the configurations were reached}
Table~\ref{tab:ablations} lists the development decisions with the largest effects. It covers 55 logged runs, all
scored on validation.

\emph{Two decisions had the largest effects.} The first is locating the word. Replacing the centred energy window
with silence trimming raised the validation log loss of A3 from 0.144 to 0.759 and of A4 from 0.119 to 0.276. For
the HMM, per-utterance CMVN and the energy window together cut the log loss from 1.65 to 0.38, before any change to
the model. The second, tested on A3, is how the sequence is summarised. A3 with its last hidden state instead of
attention fell from 96.0\% to 69.7\% accuracy, and a forward-only A3 lost 0.09 log loss.

\emph{Augmentation helped every neural model, but in different forms.} SpecAugment gave A3 its largest tuning gain
(0.214 to 0.153), but adding time shift, tempo and noise on top made it worse (0.166). For A4, SpecAugment alone changed
nothing (0.176 to 0.179), while the full set (SpecAugment, time shift, tempo and noise) improved it to 0.140, a gain
of about 1.5 of its seed standard deviations. The waveform counterpart improved A5 (0.111 to 0.101).

\emph{Pretraining needs fine-tuning.} A5 with a frozen encoder reached only 27\% accuracy. English-only
wav2vec~2.0 Base came close to the multilingual XLS-R (0.120 against 0.111), although it is 3.3 times smaller.

\emph{Some differences are within seed noise.} Retraining each final configuration with three seeds gave standard
deviations of 0.004 (A3), 0.023 (A4) and 0.001 (A5) in validation log loss. Differences smaller than these are not
interpreted.

\begin{table*}[t]
  \centering
  \caption{Development decisions with the largest effects (validation log loss, seed 42). Each row changes the
  then-best configuration; the A2 grid and A4's full augmentation change several settings at once. The seed standard
  deviation of the final configurations is 0.004 (A3), 0.023 (A4) and 0.001 (A5).}
  \label{tab:ablations}
  \small
  \begin{tabular}{@{}llrrl@{}}
    \toprule
    Approach & Change & Before & After & Decision \\
    \midrule
    A1 & RBF SVM $\to$ multinomial logistic regression & 0.783 & 0.725 & kept (final) \\
    A1 & 13 $\to$ 40 MFCCs in the statistics & 0.783 & 0.880 & rejected \\
    A2 & + per-utterance CMVN & 1.650 & 0.931 & kept \\
    A2 & silence trim $\to$ centred 1.5~s energy window & 0.931 & 0.380 & kept \\
    A2 & 5 states $\times$ 2 Gaussians $\to$ 12 $\times$ 8 (grid search) & 0.380 & 0.180 & kept (final) \\
    A3 & + SpecAugment & 0.214 & 0.153 & kept \\
    A3 & SpecAugment $\to$ full augmentation & 0.153 & 0.166 & rejected \\
    A3 & attention pooling $\to$ last hidden state & 0.153 & 0.915 & rejected \\
    A3 & bidirectional $\to$ forward-only LSTM & 0.153 & 0.246 & rejected \\
    A3 & energy window $\to$ silence trim & 0.144 & 0.759 & rejected \\
    A3 & + convolutional front-end & 0.142 & 0.134 & kept (final) \\
    A4 & + SpecAugment & 0.176 & 0.179 & rejected \\
    A4 & + full augmentation (SpecAugment, shift, tempo, noise) & 0.176 & 0.140 & kept \\
    A4 & log-mel $\to$ MFCC-40 with deltas (noise then off) & 0.140 & 0.133 & kept \\
    A4 & 1.5~s $\to$ 2.0~s window & 0.133 & 0.119 & kept (final) \\
    A4 & energy window $\to$ silence trim & 0.119 & 0.276 & rejected \\
    A4 & deeper / wider kernels / narrower network & 0.119 & 0.151--0.158 & rejected \\
    A5 & multilingual XLS-R $\to$ English-only wav2vec~2.0 Base & 0.111 & 0.120 & rejected \\
    A5 & fine-tuned $\to$ frozen encoder (head only) & 0.111 & 1.960 & rejected \\
    A5 & + waveform augmentation (shift, speed, noise) & 0.111 & 0.101 & kept (final) \\
    A5 & 1.5~s $\to$ 1.0~s window & 0.101 & 0.130 & rejected \\
    \bottomrule
  \end{tabular}
\end{table*}

\subsection{Test results}
Table~\ref{tab:test} gives the final scores on the 630 test clips. A5 is the best model on
every metric: log loss $0.093\pm0.001$, accuracy $98.1\pm0.1\%$ and macro-F1 0.982 (12 errors for seed 42). A4
follows (log loss 0.125, 96.3\%), then A3 (0.170, 96.0\%), the HMM (0.186, 96.5\%) and the order-free A1 (0.700,
79.5\%).

Validation predicted test well. A5, A2 and A1 moved by at most 0.025 log loss, and A4 improved slightly. A3 lost the
most (0.137 to 0.170). Its configuration combines five search decisions, each chosen on validation, so this is
consistent with selection optimism~\cite{cawley2010overfitting}. With 630 clips, the seed-42 bootstrap intervals
span about $\pm$1--1.5 accuracy points for the order-aware models (for A5, 97.0--99.0\%) and $\pm$3 for A1.

\begin{table*}[t]
  \centering
  \caption{Test results (630 clips; mean $\pm$ SD over three seeds for A3--A5). Params counts all weights. Train is
  the final run's training time (T4 GPU for A3--A5, CPU for A1--A2). CPU ms/clip is the latency
  from audio file to probabilities on one CPU thread.}
  \label{tab:test}
  \small
  \begin{tabular}{@{}lrrrcccc@{}}
\toprule
Approach & Params & Train (min) & CPU ms/clip & Log loss & Accuracy (\%) & Macro-F1 & ECE \\
\midrule
A1 MFCC stats + LR & 2K & 0.1 & 7 & 0.700 & 79.5 & 0.795 & 0.047 \\
A2 GMM-HMM & 93K & 11.0 & 113 & 0.186 & 96.5 & 0.965 & 0.023 \\
A3 BiLSTM + attention & 855K & 1.2 & 14 & 0.170 $\pm$ 0.009 & 96.0 $\pm$ 0.3 & 0.960 $\pm$ 0.003 & 0.017 $\pm$ 0.001 \\
A4 TC-ResNet & 150K & 1.0 & 9 & 0.125 $\pm$ 0.026 & 96.3 $\pm$ 1.2 & 0.963 $\pm$ 0.012 & 0.019 $\pm$ 0.003 \\
A5 XLS-R (fine-tuned) & 316M & 13.0 & 871 & 0.093 $\pm$ 0.001 & 98.1 $\pm$ 0.1 & 0.982 $\pm$ 0.001 & 0.014 $\pm$ 0.003 \\
\bottomrule
\end{tabular}

\end{table*}

\begin{table}[t]
  \centering
  \caption{Significance on the test split (seed 42). McNemar's exact test counts the clips where only one model is
  right; $p$ is Holm-corrected over the five comparisons. $\Delta$ is the difference in log loss (first $-$ second)
  with a paired-bootstrap 95\% interval.}
  \label{tab:significance}
  \footnotesize
  \setlength{\tabcolsep}{3pt}
  \begin{tabular}{@{}lrrrl@{}}
    \toprule
    Comparison & Only 1st & Only 2nd & $p$ & $\Delta$ log loss [95\% CI] \\
    \midrule
    A5 vs A1 & 118 & 1 & $<10^{-32}$ & $-0.608$ [$-0.698$, $-0.531$] \\
    A5 vs A2 & 12 & 2 & 0.039 & $-0.094$ [$-0.133$, $-0.055$] \\
    A5 vs A3 & 13 & 1 & 0.007 & $-0.074$ [$-0.121$, $-0.029$] \\
    A5 vs A4 & 7 & 0 & 0.039 & $-0.014$ [$-0.036$, $+0.005$] \\
    A3 vs A4 & 4 & 9 & 0.27 & $+0.059$ [$+0.018$, $+0.101$] \\
    \bottomrule
  \end{tabular}
\end{table}

\paragraph{Which differences are real}
A5 is significantly more accurate than every other model after Holm's correction (Table~\ref{tab:significance}).
On log loss, its lead over A1--A3 is clear, but its lead over A4 is not: the interval includes zero. A5 makes fewer
errors, while A4's probabilities are nearly as good. The central architectural comparison, recurrence (A3) against
convolution (A4), shows no significant difference in accuracy (4 against 9 discordant clips, $p=0.27$). A4's log
loss is lower, however, by 0.059 [0.018, 0.101]. That holds for seed 42; across seeds the gap (0.170 against 0.125)
is under two of A4's seed standard deviations. The final A3 also has a convolutional front-end, so this compares a
recurrent network with convolutional input layers against a purely convolutional one. The HMM ties with both on
accuracy as well (11 against 9 discordant clips with A3, $p=0.82$; 5 against 8 with A4, $p=0.58$; uncorrected tests
outside the family of Table~\ref{tab:significance}).

\paragraph{Calibration}
Label smoothing leaves every neural model under-confident, and the HMM's raw scores are flatter still (test ECE
0.07--0.23 before scaling). Temperatures fitted on validation (0.38--0.71) transfer to the test split and bring ECE
down to 0.011--0.023.

\paragraph{Where the errors are}
The anagram pair shows the effect of modelling order: \textit{sita}/\textit{tisa} is confused in 8.6\% of its test
clips by A1, but in only 1.0--1.9\% by A2--A5. Per word, A1 fails broadly, while the order-aware models' errors are
spread thinly over many words (Fig.~\ref{fig:perword}). Section~\ref{sec:errors} shows that ten recordings, shared by
all order-aware models, account for 10 of A5's 12 errors and about half of the other models' errors.

\begin{figure}[t]
  \centering
  \includegraphics[width=\linewidth]{figures/test_per_word_f1.png}
  \caption{Per-word test F1 (seed 42). The order-free A1 is far behind on every word. The order-aware models stay
  above 93\% on every word.}
  \label{fig:perword}
\end{figure}

\subsection{How much labelled speech is needed}
Each final configuration was retrained on 10, 25 and 50\% of the training clips (nested, stratified subsets; one
seed each) and scored on the test split (Fig.~\ref{fig:datasize}).
\begin{itemize}
  \item \emph{Pretraining pays off most when data is scarce.} With 24 clips per word, A5 reaches 96.5\%. That is about
        as accurate as the models trained from scratch on all 245 clips per word (96.2--97.0\%). It saturates by about
        120 clips per word.
  \item \emph{Without pretraining, the HMM is the most data-efficient model.} At 24 clips per word it reaches
        89.2\%, against 83.5\% for A3 and 62.5\% for A4. From 61 clips per word on, A3 and A4 stay within
        1.4 points of the HMM, and only on the full training set does A4 overtake it.
\end{itemize}
The recipe was held fixed, so the smallest subsets also received fewer optimisation steps. At 10\%, A3--A5 were
still improving at their final epoch, so those points understate them, A4 most of all.

\begin{figure*}[t]
  \centering
  \includegraphics[width=0.78\linewidth]{figures/data_size_curve.png}
  \caption{Test accuracy and log loss when each final configuration is retrained on 10--100\% of the training clips
  (24--245 per word; one seed per point).}
  \label{fig:datasize}
\end{figure*}

\subsection{Cost}
The accuracy has a price (Table~\ref{tab:test}). A5 has 316M parameters (1.26~GB as 32-bit weights) and needs
871~ms per clip on one CPU thread. A4 has 150K parameters (0.6~MB) and needs 9~ms. A5 is thus about 96 times slower and 2{,}100 times
larger, in exchange for 0.032 lower log loss and 1.8 more accuracy points. On a phone or a low-cost voice server, A4
is the efficient choice. Where an error is costly, for example when confirming a mobile-money amount, A5 is worth its
price.

\subsection{Discussion}
We return to the research question: how effectively do sequential models recognise spoken Swahili words, and why?
\begin{enumerate}
  \item \emph{Modelling order is necessary.} Every order-aware model reaches 96--98\% test accuracy, against 79.5\%
        for order-free statistics. A1 also differs in features and capacity, so the gap alone does not prove that
        order is the cause. The anagram pair and the stress test (Section~\ref{sec:errors}) make the link directly,
        as the EDA predicted: the order-aware models confuse \textit{sita} and \textit{tisa} rarely, and fail on
        reversed words. For A5, the order that matters lies mostly within about 100~ms.
  \item \emph{How order is modelled matters less than expected.} A classical HMM ties the neural models trained from
        scratch on accuracy. Convolution matches a BiLSTM with a convolutional front-end in accuracy, with 5.7 times
        fewer parameters, and has a lower log loss for seed 42. This agrees with the finding that convolutional
        networks rival recurrent ones on sequence tasks~\cite{bai2018empirical}, here for sequences as short as a
        single word.
  \item \emph{Pretraining brings the largest gain we observed,} above all with little data. A model that heard only
        91 hours of Swahili during pretraining~\cite{babu2022xlsr}, and an English-only model almost as well,
        transfer effectively to Swahili words. Our design confounds pretraining with A5's architecture, input and
        size (Section~\ref{sec:limitations}), but the result supports self-supervised pretraining as the practical
        route for low-resource languages~\cite{baevski2020wav2vec,doumbouya2021using}.
  \item \emph{The data sets the ceiling.} Locating the word in a long, noisy recording mattered as much as the
        model. Ten test clips defeat every order-aware architecture. They are quiet, noisy or badly cropped far more
        often than other clips, and include a likely label error and words said in English
        (Section~\ref{sec:errors}).
\end{enumerate}
